\documentclass[12pt]{article}

\usepackage{sbc-template}
\usepackage{graphicx,url}
\usepackage[utf8]{inputenc}
\usepackage[brazil]{babel}
\usepackage{booktabs}
\usepackage{multirow}
\usepackage{float}
\usepackage{amsmath}

\sloppy

\title{Análise Estatística de Métricas de Código-Fonte\\do Framework Neo4j}

\author{Igor Mariano Lopes Rodrigues\inst{1}, Filipe Oliveira de Saldanha da Gama\inst{1}}

\address{IBMEC -- Curso de Graduação em Engenharia de Software\\
  \email{\{202407095992,202303161085\}@alunos.ibmec.edu.br}
}

\begin{document}

\maketitle

\begin{abstract}
This paper presents a statistical analysis of source code metrics from the Neo4j
framework using the GitHub Bug Dataset v1.1, covering class-level (7{,}917 instances)
and file-level (4{,}261 instances) granularities. Descriptive statistics (central
tendency, dispersion, quartiles), normality tests (Shapiro-Wilk and Lilliefors),
Pearson and Spearman correlations, VIF multicollinearity analysis, cost-sensitive
logistic regression, and 10-fold cross-validation were applied. Strong correlations
among size and complexity metrics were found (e.g., LOC$\times$LLOC: $r = 0.99$).
All distributions are non-normal, and a severe class imbalance ($<0.15\%$ positive)
fundamentally limits predictive modeling---with only 8 buggy classes and 6 buggy
files, the Events Per Variable ratio is far below the recommended minimum of 10,
rendering any classifier unreliable. The main finding is dataset inadequacy for fault
prediction, not the inadequacy of the metrics themselves.

\noindent\textbf{Keywords:} Software metrics; Object-oriented design; Statistical analysis; Fault prediction; Class imbalance; Neo4j.
\end{abstract}

\begin{resumo}
Este artigo apresenta uma análise estatística de métricas de código-fonte do
framework Neo4j, utilizando o GitHub Bug Dataset v1.1, nos níveis classe (7.917
instâncias) e arquivo (4.261 instâncias). Foram aplicados estatística descritiva
(tendência central, dispersão, quartis), testes de normalidade (Shapiro-Wilk e
Lilliefors), correlações de Pearson e Spearman, análise VIF, regressão logística
\textit{cost-sensitive} e validação cruzada com 10 \textit{folds}. Correlações
fortes foram encontradas entre métricas de tamanho e complexidade (LOC$\times$LLOC:
$r = 0{,}99$). Todas as distribuições são não-normais e o desbalanceamento severo
($<0{,}15\%$ positivos) limita fundamentalmente a modelagem preditiva---com apenas
8 classes com bug e 6 arquivos com bug, a razão EPV está muito abaixo do mínimo
recomendado de 10, tornando qualquer classificador não confiável. O principal
achado é a inadequação do dataset para predição de defeitos, e não a inadequação
das métricas em si.

\noindent\textbf{Palavras-chave:} Métricas de software; Projeto orientado a objetos; Análise estatística; Predição de defeitos; Desbalanceamento de classes; Neo4j.
\end{resumo}

% =====================================================================
\section{Introdução} \label{sec:intro}

A análise quantitativa de métricas de código-fonte constitui abordagem fundamental na Engenharia de Software para compreender qualidade, complexidade e propensão a defeitos \cite{pressman2014}. As métricas orientadas a objetos propostas por Chidamber e Kemerer \cite{chidamber1994} permitem avaliar características estruturais de projetos Java e são amplamente utilizadas em predição de falhas \cite{gyimothy2005, basili1996}.

O framework Neo4j é um sistema de banco de dados em grafo de código aberto, amplamente adotado em aplicações que demandam modelagem de relacionamentos complexos. Por sua relevância no ecossistema Java e pelo volume considerável de código-fonte, constitui um caso de estudo relevante para análise de métricas de software.

Este trabalho aplica técnicas de estatística descritiva e inferencial sobre datasets de métricas do framework Neo4j, extraídos do GitHub Bug Dataset v1.1 \cite{toth2016}. A análise contempla: (i)~medidas de tendência central, dispersão e posição relativa; (ii)~gráficos e detecção de outliers; (iii)~testes de normalidade (Shapiro-Wilk e Lilliefors); (iv)~correlação (Pearson e Spearman) e multicolinearidade (VIF); e (v)~regressão logística \textit{cost-sensitive} com validação cruzada. O repositório está disponível em: \url{https://github.com/IgorMariano25/Projeto-ES-Neo4j}.

% =====================================================================
\section{Metodologia} \label{sec:metodo}

\subsection{Fonte de Dados}

Os datasets foram obtidos do GitHub Bug Dataset v1.1 \cite{toth2016}. Para o Neo4j, foram selecionados os datasets mais recentes: \textbf{Class-level} com 7.917 instâncias e 113 atributos (complexidade, acoplamento, coesão, herança, tamanho, documentação, warnings) e \textbf{File-level} com 4.261 instâncias e 12 atributos (complexidade ciclomática McCC, LLOC e histórico de modificações). Nenhum dataset apresentou valores ausentes.

\subsection{Seleção de Variáveis}

A seleção aplicou três critérios objetivos sobre os 113 atributos Class-level e os 12 atributos File-level. \textbf{(i) Natureza métrica}: foram retidas apenas variáveis com semântica de propriedade quantificável do código-fonte; foram descartados identificadores textuais (\texttt{ID}, \texttt{Name}, \texttt{LongName}, \texttt{Parent}, \texttt{Component}, \texttt{Path}) e metadados posicionais (\texttt{Line}, \texttt{Column}, \texttt{EndLine}, \texttt{EndColumn}), que não constituem medidas de qualidade ou de propensão a defeitos. \textbf{(ii) Variabilidade não-nula}: variáveis constantes foram excluídas, pois invalidam correlações, testes de hipótese e o cálculo do VIF; no File-level esse critério eliminou \textbf{CLOC} (\textit{Comment Lines of Code}), cuja totalidade das 4.261 instâncias é igual a zero (verificado por inspeção direta). \textbf{(iii) Não-redundância}: no Class-level, foram descartadas as variantes com prefixo \textbf{T} (totais agregados a classes aninhadas, com $r>0{,}99$ contra suas contrapartes), as variantes locais (prefixo \texttt{NL}) e proxies redundantes (CBOI, NII, NOI) já capturados por CBO/RFC; as 37 categorias granulares de regras PMD foram colapsadas nos agregados \texttt{WarningInfo} e \texttt{WarningMajor}.

Restam \textbf{19 variáveis numéricas no Class-level} (Tabela~\ref{tab:selecao}) e \textbf{7 no File-level}: McCC, LLOC, \textit{Number of previous modifications}, \textit{Number of previous fixes}, \textit{Number of committers}, \textit{Number of developer commits} e a variável-alvo \textit{Number of bugs}. O total exato de sete decorre da aritmética do dataset (12 colunas $-$ 4 identificadores $-$ 1 constante), correspondendo a \emph{todas} as métricas remanescentes após os critérios (i) e (ii), sem corte subjetivo adicional.

\begin{table}[H]
\centering
\caption{Variáveis selecionadas para análise (Class-level)}
\label{tab:selecao}
\footnotesize
\begin{tabular}{llp{7.5cm}}
\toprule
\textbf{Variável} & \textbf{Categoria} & \textbf{Descrição} \\
\midrule
WMC, NL, NLE             & Complexidade  & WMC (\textit{Weighted Methods per Class}), NL (\textit{Nesting Level}), NLE (\textit{Nesting Level Else-if}) \\
CBO, RFC                 & Acoplamento   & CBO (\textit{Coupling Between Objects}), RFC (\textit{Response For a Class}) \\
LCOM5                    & Coesão        & LCOM5 (\textit{Lack of Cohesion in Methods}, versão 5) \\
DIT, NOC                 & Herança       & DIT (\textit{Depth of Inheritance Tree}), NOC (\textit{Number of Children}) \\
LOC, LLOC, NOS, NM, NA   & Tamanho       & LOC (\textit{Lines of Code}), LLOC (\textit{Logical Lines of Code}), NOS (\textit{Number of Statements}), NM (\textit{Number of Methods}), NA (\textit{Number of Attributes}) \\
CC                       & Clones        & CC (\textit{Clone Coverage}) --- proporção de código em clones \\
CD, CLOC                 & Documentação  & CD (\textit{Comment Density}), CLOC (\textit{Comment Lines of Code}) \\
WarningInfo, WarningMajor & Qualidade    & Warnings PMD agregados por severidade: \textit{Info} (informativos) e \textit{Major} (majoritários) \\
N.\ de bugs              & Alvo          & \textit{Number of bugs} --- número de defeitos registrados \\
\bottomrule
\end{tabular}
\end{table}

% =====================================================================
\section{Resultados e Discussão} \label{sec:resultados}

\subsection{Medidas de Tendência Central}

A Tabela~\ref{tab:tendencia_class} apresenta média, mediana e moda para as principais variáveis do dataset Class-level.

\begin{table}[H]
\centering
\caption{Medidas de tendência central --- Class-level}
\label{tab:tendencia_class}
\footnotesize
\begin{tabular}{lrrr}
\toprule
\textbf{Variável} & \textbf{Média} & \textbf{Mediana} & \textbf{Moda} \\
\midrule
WMC   &  7,39  &  3,00  &  1 \\
CBO   &  5,64  &  3,00  &  2 \\
RFC   &  9,84  &  5,00  &  1 \\
LCOM5 &  1,58  &  1,00  &  1 \\
LOC   & 59,70  & 25,00  &  7 \\
LLOC  & 49,18  & 21,00  &  8 \\
NOS   & 18,77  &  5,00  &  1 \\
NM    &  8,61  &  4,00  &  1 \\
DIT   &  0,98  &  1,00  &  1 \\
CC    &  0,12  &  0,00  &  0 \\
\bottomrule
\end{tabular}
\end{table}

Em todas as métricas de tamanho e complexidade, a média é substancialmente superior à mediana, indicando distribuições fortemente assimétricas à direita. LOC apresenta média de 59,70 e mediana de apenas 25,00: a maioria das classes possui código reduzido enquanto um subconjunto menor concentra valores muito elevados. A moda igual a 1 para WMC, RFC, NOS e NM indica que a maioria das classes possui complexidade mínima. DIT apresenta baixa variabilidade (mediana = moda = 1), refletindo hierarquias de herança rasas no Neo4j.

Para o File-level, McCC apresenta média de 6,06 e mediana de 2,00 (moda = 1), enquanto LLOC possui média de 102,74 e mediana de 55,00, confirmando o mesmo padrão de assimetria.

\subsection{Medidas de Dispersão}

A Tabela~\ref{tab:dispersao_class} apresenta amplitude, variância e desvio padrão para as principais variáveis.

\begin{table}[H]
\centering
\caption{Medidas de dispersão --- Class-level}
\label{tab:dispersao_class}
\footnotesize
\begin{tabular}{lrrr}
\toprule
\textbf{Variável} & \textbf{Amplitude} & \textbf{Variância} & \textbf{Desvio Padrão} \\
\midrule
WMC   &     325 &      193,25 &  13,90 \\
CBO   &     161 &       55,57 &   7,45 \\
RFC   &     293 &      256,19 &  16,01 \\
LCOM5 &      71 &       12,33 &   3,51 \\
LOC   &   3.207 &   13.360,81 & 115,59 \\
LLOC  &   2.588 &    8.473,03 &  92,05 \\
NOS   &   1.411 &    2.218,10 &  47,10 \\
NM    &     161 &      183,94 &  13,56 \\
DIT   &       6 &        0,99 &   0,99 \\
CC    &       1 &        0,08 &   0,29 \\
\bottomrule
\end{tabular}
\end{table}

LOC apresenta amplitude de 3.207 e desvio padrão de 115,59 --- quase o dobro da média (59,70) --- confirmando a alta heterogeneidade do tamanho das classes. O coeficiente de variação de LOC é de 1,94, o maior entre as métricas de tamanho. DIT, por outro lado, apresenta variância de apenas 0,99, tornando essa métrica praticamente irrelevante como discriminador de \textit{fault-proneness} no Neo4j. No File-level, ``Number of developer commits'' apresenta variância de 4.106.050,59 e desvio padrão de 2.026,34, decorrente da grande discrepância entre desenvolvedores.

\subsection{Medidas de Posição Relativa e Boxplots}

A análise de quartis complementa as medidas anteriores. Para WMC no Class-level: $Q_1 = 1$, $Q_2 = 3$, $Q_3 = 8$, IQR = 7 e $P_{95} = 27$. Isso indica que 75\% das classes possuem WMC $\leq 8$, enquanto o valor máximo é 325. Para LOC: $Q_1 = 7$, $Q_2 = 25$, $Q_3 = 60$, IQR = 53 --- 75\% das classes possuem menos de 60 linhas, mas o valor máximo é 3.207.

\begin{figure}[H]
\centering
\includegraphics[width=0.97\textwidth]{Ap1/figuras/boxplots_class.png}
\caption{Boxplots das principais métricas --- Class-level}
\label{fig:boxplots_class}
\end{figure}

A Figura~\ref{fig:boxplots_class} evidencia a presença massiva de outliers em todas as métricas, padrão típico em projetos de software de grande porte. O limite inferior de detecção ($Q_1 - 1{,}5 \times \text{IQR}$) é negativo para todas as variáveis, de modo que apenas outliers superiores foram detectados. No File-level, é observado padrão análogo: McCC ($Q_1 = 1$, $Q_3 = 6$, $P_{90} = 18$) e LLOC ($Q_1 = 20$, $Q_3 = 112$) também exibem distribuições assimétricas com outliers extremos.

\subsection{Avaliação de Outliers}

\begin{table}[H]
\centering
\caption{Detecção de outliers pelo método IQR --- Class-level}
\label{tab:outliers}
\footnotesize
\begin{tabular}{lrrr}
\toprule
\textbf{Variável} & \textbf{Outliers} & \textbf{\% Total} & \textbf{Limite Sup.} \\
\midrule
WMC   &   706 &  8,92\% &  18,50 \\
CBO   &   644 &  8,13\% &  14,50 \\
RFC   &   759 &  9,59\% &  24,50 \\
LOC   &   799 & 10,09\% & 140,00 \\
LLOC  &   787 &  9,94\% & 112,50 \\
NOS   &   840 & 10,61\% &  42,00 \\
NM    &   866 & 10,94\% &  19,50 \\
CC    & 1.780 & 22,48\% &   0,00 \\
\bottomrule
\end{tabular}
\end{table}

A proporção de outliers varia entre 5,86\% (NA) e 22,48\% (CC). Para as métricas de tamanho (LOC, LLOC, NOS), aproximadamente 10\% são outliers, consistente com a regra de Pareto em software: uma minoria de componentes concentra a maior complexidade \cite{nagappan2006}. CC apresenta o maior percentual de outliers porque a maioria dos valores é zero; qualquer valor positivo excede $Q_3 + 1{,}5 \times \text{IQR} = 0{,}00$. Os outliers não foram removidos por representarem classes legitimamente complexas do framework.

\subsection{Visualizações Gráficas}

A inspeção visual das distribuições complementa as estatísticas descritivas apresentadas nas seções anteriores, evidenciando padrões de forma, simetria e relacionamento entre variáveis que medidas resumo isoladas não revelam. Esta seção apresenta os histogramas das principais métricas no nível de classe, que permitem avaliar a forma das distribuições marginais, e um gráfico de dispersão bivariada que ilustra a estrutura conjunta entre tamanho (LOC) e complexidade (WMC) --- duas dimensões fortemente correlacionadas no Neo4j.

\begin{figure}[H]
\centering
\includegraphics[width=0.97\textwidth]{Ap1/figuras/histogramas_class.png}
\caption{Histogramas das métricas Class-level}
\label{fig:hist_class}
\end{figure}

Os histogramas (Figura~\ref{fig:hist_class}) revelam \textbf{assimetria positiva (à direita)} em \emph{todas} as oito métricas representadas (WMC, CBO, RFC, LOC, LLOC, NOS, NM e LCOM5): grande concentração de massa em valores baixos e caudas longas estendendo-se à direita, sem qualquer ocorrência de assimetria à esquerda. A intensidade do enviesamento varia em grau: LOC, LLOC e NOS exibem as caudas mais longas (limites superiores observados de 3.207, 2.588 e 1.411, respectivamente), enquanto WMC, RFC e LCOM5 mostram pico extremo próximo de zero, com mais de 50\% das observações concentradas no primeiro \textit{bin} do histograma.

Essa forma carrega três implicações estatísticas relevantes. \textbf{(a) Tendência central distorcida.} A relação $\bar{x} > \text{mediana} > \text{moda}$, já evidente na Tabela~\ref{tab:tendencia_class} (e.g., LOC: 59,70 vs.\ 25 vs.\ 7; WMC: 7,39 vs.\ 3 vs.\ 1), é o marcador algébrico clássico de assimetria positiva: a média é inflada pelas caudas e não representa o comportamento típico das classes, de modo que mediana e quartis tornam-se descritores preferíveis. \textbf{(b) Inadequação de métodos paramétricos.} Técnicas que pressupõem normalidade --- teste $t$, ANOVA, intervalos de confiança gaussianos e regressão linear sem transformação --- tornam-se enviesadas neste contexto, o que justifica adiante o uso da correlação de Spearman (baseada em postos) em complemento a Pearson e a opção por regressão logística com pesos \textit{cost-sensitive} em vez de regressão linear. \textbf{(c) Compatibilidade com distribuições típicas de software.} A combinação pico-em-zero, decaimento aproximadamente exponencial e cauda pesada é compatível com distribuições \textit{log-normal} ou de \textit{lei de potência}, padrão amplamente reportado para métricas de tamanho e acoplamento em sistemas de grande porte~\cite{radjenovic2013}.

A inspeção visual converge com os testes inferenciais discutidos na Seção~\ref{sec:normalidade}: distribuições visualmente não-simétricas e sem o formato \textit{bell-shaped} característico da gaussiana fazem a estatística $W$ de Shapiro-Wilk afastar-se de 1 e a estatística $D$ de Lilliefors crescer, com $p < 10^{-50}$ uniformemente em todas as variáveis. Gráfico e teste contam a mesma história por meios diferentes, e essa convergência reforça a robustez do diagnóstico de não-normalidade.

A distribuição de bugs, por sua vez, é extremamente desbalanceada: apenas 8 classes com bug (0,10\%) e 6 arquivos com bug (0,14\%) entre milhares de instâncias --- desbalanceamento que constituirá o principal fator limitante da modelagem preditiva discutida adiante.

\begin{figure}[H]
\centering
\includegraphics[width=0.62\textwidth]{Ap1/figuras/dispersao_loc_wmc.png}
\caption{Gráfico de dispersão: LOC vs WMC (Class-level)}
\label{fig:disp_loc_wmc}
\end{figure}

A Figura~\ref{fig:disp_loc_wmc} apresenta a dispersão entre LOC e WMC, evidenciando a correlação forte e positiva entre tamanho e complexidade. As poucas instâncias com bugs (pontos vermelhos) tendem a se posicionar em regiões de valores moderados a altos, embora o número ínfimo de positivos (8 classes com bug) não permita conclusões robustas sobre a relação com defeitos.

\subsection{Testes de Normalidade} \label{sec:normalidade}

Para todas as variáveis, os testes de Shapiro-Wilk \cite{shapiro1965} e Lilliefors \cite{lilliefors1967} rejeitaram a hipótese nula de normalidade ($p < 10^{-50}$). A Tabela~\ref{tab:normalidade} apresenta as estatísticas para variáveis representativas.

\begin{table}[H]
\centering
\caption{Testes de normalidade --- Class-level (amostra de 5.000 obs.\ para SW; seed=42)}
\label{tab:normalidade}
\footnotesize
\begin{tabular}{lcccc}
\toprule
\textbf{Variável} & \textbf{SW} & \textbf{SW $p$} & \textbf{Lill.\ $D$} & \textbf{Lill.\ $p$} \\
\midrule
WMC & 0,4214 & $< 10^{-83}$   & 0,3081 & $< 2{,}2\!\times\!10^{-16}$ \\
CBO & 0,5874 & $< 10^{-76}$   & 0,2246 & $< 2{,}2\!\times\!10^{-16}$ \\
RFC & 0,5158 & $< 10^{-79}$   & 0,2786 & $< 2{,}2\!\times\!10^{-16}$ \\
LOC & 0,4244 & $< 10^{-83}$   & 0,3113 & $< 2{,}2\!\times\!10^{-16}$ \\
DIT & 0,8267 & $< 10^{-59}$   & 0,2520 & $< 2{,}2\!\times\!10^{-16}$ \\
CC  & 0,4754 & $< 10^{-81}$   & 0,4410 & $< 2{,}2\!\times\!10^{-16}$ \\
\bottomrule
\end{tabular}
\end{table}

A estatística de Shapiro-Wilk varia entre 0,013 (Number of bugs) e 0,827 (DIT), sempre muito distante de 1 (normalidade perfeita).

A escolha pelo teste de Lilliefors \cite{lilliefors1967} em detrimento do tradicional Kolmogorov-Smirnov (K-S) decorre de uma restrição técnica do K-S clássico: sua tabela de valores críticos só é corretamente calibrada quando a média ($\mu$) e o desvio padrão ($\sigma$) da distribuição teórica de referência são fixados \textit{a priori}, antes da observação dos dados. No presente estudo, $\mu$ e $\sigma$ não são valores populacionais conhecidos para métricas como LOC, WMC ou LLOC --- precisam ser estimados a partir da própria amostra. Sob essa estimação, o K-S aplicado com a tabela padrão torna-se \textbf{excessivamente conservador}: como os parâmetros são ajustados aos dados, a distância máxima $D$ entre a função de distribuição empírica e a teórica fica artificialmente reduzida, fazendo o teste \emph{deixar de rejeitar} a hipótese nula de normalidade com mais frequência do que o nível $\alpha$ nominal indicaria. Em linguagem acessível: nessa configuração, o K-S tende a ``aceitar como normal'' distribuições que de fato não são, justamente porque a referência teórica foi ajustada para se parecer com a amostra. O teste de Lilliefors corrige essa distorção recalculando os valores críticos de modo a respeitar o nível de significância nominal mesmo quando $\mu$ e $\sigma$ são estimados amostralmente, restaurando a calibração correta da taxa de erro tipo I. Por essa razão, ambos os testes (Shapiro-Wilk e Lilliefors) foram reportados em conjunto: a convergência dos dois testes --- com $p < 10^{-50}$ em todas as variáveis --- elimina qualquer dúvida sobre a violação de normalidade.

Para o Shapiro-Wilk, cuja implementação em R impõe limite de 5.000 observações, foi utilizada uma amostra aleatória extraída com \texttt{seed=42}, garantindo reprodutibilidade. A rejeição universal de normalidade justifica o uso da correlação de Spearman para análise de associações e impõe cautela quanto a métodos paramétricos que assumem distribuição normal.

\subsection{Análise de Correlação}

Dado que todas as variáveis violam normalidade, a correlação de Pearson foi complementada com a de Spearman (baseada em postos), mais robusta a outliers e distribuições assimétricas.

\begin{figure}[H]
\centering
\includegraphics[width=0.80\textwidth]{Ap1/figuras/correlacao_heatmap_class.png}
\caption{Correlação de Pearson --- Class-level}
\label{fig:corr_class}
\end{figure}

A Figura~\ref{fig:corr_class} apresenta a matriz de correlação de Pearson. \textbf{Correlações fortes} ($|r| \geq 0{,}7$, Pearson): LOC$\times$LLOC: $r = 0{,}99$; LOC$\times$NOS: $r = 0{,}95$; WMC$\times$LOC: $r = 0{,}89$; WMC$\times$RFC: $r = 0{,}82$; CBO$\times$RFC: $r = 0{,}80$. \textbf{Moderadas} ($0{,}4 \leq |r| < 0{,}7$): WMC$\times$CBO: $r = 0{,}51$; WMC$\times$NM: $r = 0{,}53$; CBO$\times$LOC: $r = 0{,}59$.

As correlações de Spearman (Figura~\ref{fig:corr_spearman_class}) apresentam padrão semelhante com magnitudes geralmente mais elevadas (e.g., LOC$\times$LLOC: $\rho = 0{,}99$; NOS$\times$NL: $\rho = 0{,}72$). A convergência de ambos os métodos reforça a robustez dos achados. Pearson é sensível a outliers e pressupõe linearidade; Spearman captura relações monotônicas sem essa restrição.

\begin{figure}[H]
\centering
\includegraphics[width=0.80\textwidth]{Ap1/figuras/correlacao_spearman_class.png}
\caption{Correlação de Spearman --- Class-level}
\label{fig:corr_spearman_class}
\end{figure}

\begin{figure}[H]
\centering
\includegraphics[width=0.80\textwidth]{Ap1/figuras/chart_correlation_class.png}
\caption{Correlação de Pearson com dispersões e histogramas marginais (\texttt{chart.Correlation}) --- Class-level}
\label{fig:chart_corr}
\end{figure}

A Figura~\ref{fig:chart_corr}, gerada com o pacote \texttt{PerformanceAnalytics} sobre o subconjunto Class-level, combina histogramas (diagonal), gráficos de dispersão (triângulo inferior) e coeficientes de correlação com nível de significância (triângulo superior), proporcionando uma visão integrada das relações entre as variáveis no nível de classe. Para o dataset File-level (figura análoga não reproduzida aqui por concisão), as correlações mais fortes ocorrem entre McCC$\times$LLOC ($r = 0{,}72$) e ``N.\ previous modifications''$\times$``N.\ committers'' ($r = 0{,}77$; Spearman: $\rho = 0{,}90$).

\subsection{Análise de Multicolinearidade (VIF)}

O \textit{Variance Inflation Factor} (VIF) foi calculado para quantificar formalmente a multicolinearidade entre os preditores \cite{montgomery2014}. VIF $> 10$ indica multicolinearidade severa.

\begin{figure}[H]
\centering
\includegraphics[width=0.85\textwidth]{Ap1/figuras/vif_barplot.png}
\caption{VIF dos preditores da regressão logística (antes da remoção)}
\label{fig:vif}
\end{figure}

No Class-level, RFC (VIF = 20,18), WMC (15,17) e LOC (10,98) excederam o limiar de 10 (Figura~\ref{fig:vif}). Após remoção iterativa de RFC, os VIF finais ficaram abaixo de 10 (WMC: 8,75; LOC: 9,16; demais $\leq$ 5). No File-level, todos os VIF foram aceitáveis ($\leq$ 5,58), indicando ausência de multicolinearidade severa. A remoção de RFC do modelo Class-level é justificada pela sua redundância com WMC ($r = 0{,}82$) e CBO ($r = 0{,}80$).

\subsection{Modelagem Estatística --- Regressão Logística}

Foi adotada a regressão logística binária com pesos \textit{cost-sensitive} inversamente proporcionais à frequência de cada classe ($w_+ = N / (2 \cdot N_+)$; $w_- = N / (2 \cdot N_-)$), evitando classificação trivial como classe majoritária.

\subsubsection{Modelo Class-level}

O modelo utilizou os preditores remanescentes após remoção de VIF $> 10$: WMC, CBO, LOC, LCOM5, NL, DIT, CC. Os coeficientes apresentaram magnitude extrema (ordem $10^{12}$--$10^{15}$), indicando que o algoritmo \textbf{não convergiu} --- característica de separação completa dos dados, sintoma do desbalanceamento extremo (8 classes com bug em 7.917, EPV $\approx 1$). O Pseudo $R^2$ (McFadden) foi $-10{,}76$ (modelo pior que o nulo). Na escala de McFadden, valores entre 0,2 e 0,4 indicam ajuste excelente \cite{mcfadden1974}.

\subsubsection{Modelo File-level}

Os preditores foram: McCC, LLOC, N.\ previous modifications, N.\ previous fixes, N.\ committers e N.\ developer commits. O Pseudo $R^2$ (McFadden) foi 0,59 (excelente na escala), porém o algoritmo também não convergiu. Apenas LLOC não foi significativo ($p = 0{,}30$); os demais resultados devem ser interpretados com cautela dado o número ínfimo de positivos (6 arquivos com bug).

\subsubsection{Métricas de Avaliação e Validação Cruzada}

Para evitar overfitting (o modelo ajustado e avaliado no mesmo conjunto gera estimativas otimistas), foi implementada validação cruzada estratificada com 10 \textit{folds}, garantindo que cada \textit{fold} contenha pelo menos uma instância positiva. O threshold ótimo foi determinado pelo Youden Index ($J = \text{Sensit.} + \text{Espec.} - 1$).

\begin{table}[H]
\centering
\caption{Métricas de avaliação: treino completo e 10-fold CV}
\label{tab:metricas_avaliacao}
\footnotesize
\begin{tabular}{lrrrr}
\toprule
 & \multicolumn{2}{c}{\textbf{Treino completo}} & \multicolumn{2}{c}{\textbf{10-fold CV}} \\
\cmidrule(lr){2-3} \cmidrule(lr){4-5}
\textbf{Métrica} & \textbf{Class} & \textbf{File} & \textbf{Class} & \textbf{File} \\
\midrule
AUC-ROC     & 0,674 & 0,968 & 0,660 & 0,805 \\
AUC-PR      & 0,007 & 0,020 & ---   & ---   \\
Sensitivity & 0,375 & 1,000 & 0,625 & 0,833 \\
Precision   & 0,014 & 0,024 & 0,002 & 0,007 \\
F1-score    & 0,027 & 0,047 & 0,004 & 0,014 \\
Bal.\ Acc.  & 0,674 & 0,972 & 0,664 & 0,836 \\
MCC         & 0,069 & 0,151 & 0,023 & 0,068 \\
\bottomrule
\end{tabular}
\end{table}

\begin{figure}[H]
\centering
\includegraphics[width=0.97\textwidth]{Ap1/figuras/curva_roc.png}
\caption{Curvas ROC com Youden Index para Class-level e File-level}
\label{fig:roc}
\end{figure}

A Figura~\ref{fig:roc} apresenta as curvas ROC com os pontos ótimos pelo Youden Index. A degradação na validação cruzada confirma overfitting, especialmente no File-level (AUC: 0,97 $\to$ 0,81). Os valores de MCC próximos de zero (0,023 e 0,068 no CV) indicam poder preditivo praticamente nulo. A ``acurácia'' de 97--99\% é enganosa: um classificador trivial que sempre prediz ``sem bug'' obtém desempenho similar. As curvas Precision-Recall (AUC-PR $\approx 0{,}007$ e $0{,}020$) confirmam a inadequação dos dados para predição, pois uma base de referência aleatória teria AUC-PR igual à prevalência da classe positiva ($\approx 0{,}001$).

\subsection{Discussão}

As métricas do Neo4j seguem distribuições fortemente assimétricas, consistente com a literatura \cite{radjenovic2013, gyimothy2005}. As correlações fortes entre LOC, LLOC, NOS, WMC e RFC indicam dimensões sobrepostas de tamanho e complexidade, confirmadas pela análise VIF. A correlação extremamente fraca entre métricas e bugs (máximo $r = 0{,}06$) contrasta com estudos anteriores \cite{basili1996, gyimothy2005}, explicável pelo desbalanceamento extremo e provável subreporte de defeitos.

Quanto às métricas CK: DIT e NOC possuem variância quase zero no Neo4j (design ``flat''), tornando-as pouco discriminativas. LCOM5 possui múltiplas definições na literatura, dificultando comparações entre estudos. As métricas WMC e CBO, apontadas como melhores preditores na literatura \cite{radjenovic2013}, também não apresentaram poder preditivo aqui --- o que reforça que a limitação é do dataset, não das métricas.

% =====================================================================
\subsection{Ameaças à Validade} \label{sec:ameacas}

\textbf{Validade interna.} (i) Desbalanceamento extremo (8 classes com bug e 6 arquivos com bug; EPV $\approx 1$, muito abaixo do mínimo de 10 \cite{hosmer2013}), responsável pela separação completa e pela não-convergência da regressão logística; (ii) K-S com parâmetros estimados seria excessivamente conservador --- corrigido com Lilliefors; (iii) Multicolinearidade severa (LOC$\times$LLOC: $r = 0{,}99$) --- corrigida via remoção iterativa de variáveis com VIF $> 10$; (iv) Modelo originalmente avaliado no próprio conjunto de treino --- corrigido com validação cruzada estratificada de 10 \textit{folds}; (v) Subamostragem aleatória para o Shapiro-Wilk (limite de 5.000 observações) com semente fixa (\texttt{seed=42}) torna a estatística reprodutível mas potencialmente sensível à amostra particular extraída --- mitigado pela convergência com Lilliefors, computado sobre a totalidade das observações.

\textbf{Validade externa.} (i) Resultados restritos a um único projeto (Neo4j) em uma única versão, sem garantia de generalização para outros sistemas Java; (ii) domínio específico (banco de dados em grafo) com arquitetura estruturalmente plana --- a observação de que DIT e NOC apresentam variância quase nula é provavelmente idiossincrasia do Neo4j e não se reproduz em projetos com hierarquias de herança mais profundas; (iii) viés de captura do GitHub Bug Dataset~\cite{toth2016}: apenas defeitos vinculados a \textit{commits} de correção são registrados, deixando bugs latentes ou corrigidos sem rastreabilidade explícita ausentes da variável-alvo.

\textbf{Validade de construto.} (i) Métricas CK são \textit{proxies} estruturais imperfeitos da qualidade real do software --- não capturam manutenibilidade, legibilidade ou aderência a contratos; (ii) ``\textit{Number of bugs}'' operacionaliza ``defeito'' como ``bug historicamente registrado e vinculado a \textit{commit}'', construto que pode divergir substancialmente da noção teórica de ``defeito existente''; (iii) LCOM5 possui múltiplas formulações na literatura, dificultando comparações entre estudos; (iv) a agregação de 37 categorias granulares de regras PMD nos dois rótulos \texttt{WarningInfo} e \texttt{WarningMajor} sacrifica resolução semântica por parcimônia analítica.

\textbf{Validade de conclusão.} (i) Baixíssimo poder estatístico para inferências sobre a variável-alvo: com apenas 8 classes com bug e 6 arquivos com bug, qualquer teste envolvendo bugs tem alta probabilidade de erro tipo II (falhar em detectar associações reais); (ii) cerca de 170 correlações pareadas foram calculadas no Class-level sem correção formal para comparações múltiplas (Bonferroni ou FDR), elevando o risco teórico de falsos positivos entre correlações marginais --- mitigação parcial: as associações reportadas como fortes ($|r| \geq 0{,}7$) coincidem em Pearson e Spearman, conferindo robustez ao diagnóstico; (iii) os $p$-valores extremamente pequenos nos testes de normalidade ($p < 10^{-50}$) refletem em parte a sensibilidade de Shapiro-Wilk e Lilliefors a desvios mínimos quando $n$ é grande --- a leitura correta é que ``normalidade é praticamente impossível dado o tamanho amostral'', e não que o desvio absoluto seja extraordinário em magnitude prática; (iv) a confiabilidade de medição da ferramenta SourceMeter é alta por ser determinística, mas suas convenções internas (escopo de classe, contagem de linhas vazias, etc.) afetam todas as métricas de modo sistemático e não são reveladas por nenhum teste estatístico.

% =====================================================================
\section{Conclusão} \label{sec:conclusao}

Este trabalho apresentou uma análise estatística rigorosa de métricas de código-fonte do framework Neo4j (GitHub Bug Dataset v1.1), com testes de normalidade corrigidos (Lilliefors), correlações de Pearson e Spearman, análise VIF, regressão \textit{cost-sensitive} e validação cruzada.

Os principais achados foram: (1) todas as métricas seguem distribuições assimétricas com correlações fortes entre tamanho e complexidade (LOC$\times$LLOC: $r = 0{,}99$); (2) nenhuma variável é normal (todos os $p < 10^{-50}$); (3) a proporção de outliers varia entre 5\% e 22\%, representando componentes legitimamente complexos; (4) o \textbf{principal achado é a inadequação do dataset para predição de defeitos}: com EPV $\approx 1$, os modelos são não confiáveis e a ``alta acurácia'' é enganosa --- o MCC próximo de zero na CV (0,023 e 0,068) confirma poder preditivo nulo.

Como trabalhos futuros: (i) combinar múltiplas versões do Neo4j para aumentar instâncias positivas e viabilizar EPV $\geq 10$; (ii) aplicar \textit{Random Forest} ou \textit{gradient boosting}, que tratam desbalanceamento nativamente; (iii) explorar \textit{cross-project defect prediction}.

% =====================================================================
\bibliographystyle{plain}
\bibliography{sbc-template}

\end{document}
